Bei konstantem Druck verrichtet das Gas die Arbeit
\begin{align}
 W_\text{Gas} &= \WNF = \pN\cdot\dV = \WN \approx \result{\WNP{0}{3}}
\end{align}
Stickstoff (N$_2$) ist zweiatomig, $f = \fNO$. Mit $p\,V = nRT$ ist bei konstantem Druck $nR\,\Delta T = p\,\Delta V$, also nimmt die innere Energie zu um
\begin{align}
 \Delta U &= \frac{f}{2}\,nR\,\Delta T = \dUNF = \frac{\fNO}{2}\cdot\pN\cdot\dV = \dUN \approx \result{\dUNP{0}{3}}
\end{align}
Nach dem ersten Hauptsatz muss die zugeführte Wärme die Zunahme der inneren Energie und die vom Gas verrichtete Arbeit decken:
\begin{align}
 Q &= \Delta U + W_\text{Gas} = \QNF \\
   &= \dUN+\WN = \QN \approx \result{\QNP{0}{3}}
\end{align}
